\subsection{外代数的正向极限}

设 \((\mathbf{A}_{\alpha},\phi_{\beta \alpha})\) 为交换环的一个有向正向系统， \((\mathbf{M}_{\alpha},f_{\beta \alpha})\) 为 \(\mathbf{A}_{\alpha}\) -模的一个正向系统， \(\mathbf{A} = \varinjlim \mathbf{A}_{\alpha}\) 与 \(\mathbf{M} = \varinjlim \mathbf{M}_{\alpha}\) 。对于 \(\alpha \leqslant \beta\) ，我们由 \(\mathbf{A}_{\alpha}\) -同态 \(f_{\beta \alpha}: \overrightarrow{\mathbf{M}_{\alpha}} \to \mathbf{M}_{\beta}\) 典范地导出一个 \(\mathbf{A}_{\alpha}\) -代数同态（第 5 号，公式 (12)） \(f_{\beta \alpha}^{\prime \prime}: \bigwedge_{\mathbf{A}_{\alpha}} (\mathbf{M}_{\alpha}) \to \bigwedge_{\mathbf{A}_{\beta}} (\mathbf{M}_{\beta})\) ，并且由 (13) 可知 \((\bigwedge_{\mathbf{A}_{\alpha}} (\mathbf{M}_{\alpha}), f_{\beta \alpha}^{\prime \prime})\) 是分次 \(\mathbf{A}_{\alpha}\) -模的一个正向系统。另一方面，设 \(f_{\alpha}: \mathbf{M}_{\alpha} \to \mathbf{M}\) 为典范 \(\mathbf{A}_{\alpha}\) -同态；我们（由第 5 号，公式 (12)）导出一个分次 \(\mathbf{A}_{\alpha}\) -代数同态 \(f_{\alpha}^{\prime \prime}: \bigwedge_{\mathbf{A}_{\alpha}} (\mathbf{M}_{\alpha}) \to \bigwedge_{\mathbf{A}} (\mathbf{M})\) ，并且由 (13) 同样得出 \(f_{\alpha}^{\prime \prime}\) 构成 \(\mathbf{A}_{\alpha}\) -同态的一个正向系统。

命题9. A-同态 \(f'' = \varinjlim f_{\alpha}' : \varinjlim \bigwedge_{A_{\alpha}}(M_{\alpha}) \to \bigwedge_{A}(M)\) 是一个分次代数同构。

证明与 §5 第 4 号命题 6 的证明相同，只需将其中所有的 T 换为 \(\bigwedge\) ，将 \(\phi\) 换为 \(\phi''\) ，并注意到交错代数的正向极限仍是交错的。

